% !TEX root = ../main.tex

\section{Introduction}
\label{sec:intro}

Historically, standard asset pricing models conclude that any risk-averse person should invest at least a fraction of their savings in the stock market to build wealth and earn the equity premium. Yet, in reality, a large portion of the population completely avoids the stock market. This phenomenon is known as the "stock market participation puzzle". 

Understanding why people choose to or not to invest is critical. Theoretically, stock markets act as a vehicle of wealth allocation across agents and across time, facilitating risk-sharing and consumption smoothing to improve allocative efficiency. It is also closely relevant for the layperson, as investing in the stock market is a primary engine of wealth creation and often a major driver of wealth inequality. 

While researchers have identified individual factors like wealth, financial literacy, and social trust that influence this decision \citep{guiso2008, vanrooij2011, bu2022}, an unresolved phenomenon remains: there are persistent geographic differences in participation across U.S. states that cannot be explained by income alone \citep{chien2017}. \citet{chien2017} show that across different income brackets, Connecticut has the highest stock market participation (SMP) rate while Mississippi has the lowest. That is, there is some "black box" driver of SMP that is related to state of residence, but the specific mechanism is not identified.

The main hypothesis in this paper is that the "black box" of geographic differences in stock market participation is the partisan divide between states The motivation comes from the intersection of two strands of literature: (1) the household finance literature on the determinants of stock market participation, and (2) the political economics literature on how political environments shape economic behavior.

Using household-level data from the Panel Study of Income Dynamics (PSID) and state-level political orientation derived from Federal Election Commission (FEC) campaign contributions, this paper explores how living in a predominantly "red" versus "blue" state affects the likelihood of an individual investing in the market.  The findings reveal a textbook case of Simpson's paradox within U.S. financial behavior. On an individual level, shifting toward a Republican leaning increases a person's likelihood of investing, as documented in prior literature. However, in the aggregate, Democratic-leaning states exhibit a significantly higher baseline rate of stock market participation. Furthermore, this geographic divide strictly influences the extensive margin (whether to enter the stock market or not), rather than the intensive margin (how much to allocate to stocks once entered). This suggests that the state-level political environment acts as a structural barrier or fixed cost to market entry, rather than altering investors' risk preferences.

This paper primarily contributes to three specific strands of literature. First, it advances the extensive literature on the stock market participation puzzle and the geographic determinants of portfolio choice. While prior studies have identified that localized factors influence investment decisions, this paper provides a novel, structural explanation for the persistent state-level variations in participation. Crucially, by showing that political geography exclusively affects the extensive margin rather than the intensive margin, it clarifies the mechanism through which geographic environments constrain financial behavior.  

Second, it adds to the growing field of partisan influence on financial decision-making. While existing research demonstrates that individual investors dynamically adjust their risk-taking based on short-term political alignment and partisan optimism, this study establishes that long-term, state-level political environments act as a structural baseline determinant of market entry.  

Finally, this paper provides a distinct empirical contribution by documenting a striking Simpson's paradox within U.S. financial behavior. By disentangling state-level aggregates from individual-level preferences, the findings warn against ecological fallacies in household finance, demonstrating that while individual Republican-leaning preferences may correlate with higher market participation, the overarching structural environment of Democratic-leaning states yields a significantly higher baseline rate of stock ownership.


The remainder of the paper is organized as follows. Section 2 reviews the related literature regarding the stock market participation puzzle and partisan influences on financial decision-making. Section 3 describes the data sources, primarily the Panel Study of Income Dynamics (PSID) and the Federal Election Commission (FEC) campaign finance data, and outlines the methodology for variable construction. Section 4 presents the main empirical results, examining the effect of state-level political environments on both the extensive margin of market participation and the intensive margin of equity shares, alongside several robustness checks. Section 5 reconciles the aggregate state-level findings with established individual-level behavioral theories, formalizing the presence of a Simpson's paradox and addressing the influence of macroeconomic timing. Section 6 discusses potential theoretical channels and structural factors that might drive these observed geographic disparities. Finally, Section 7 concludes.